\dish{醤油ラーメン}
\aussidish{Shōyu ramen}
\altdish{Ramen Tokyo style}
\pour{4}
\prep{15~min + cuisson (5~min) [plus préparation
  \hyperref[chashu]{叉焼}]}
\source{cjm}


\stage{Soupe}

\begin{ingredients}
  \ingr{1,6}{L}{ \hyperref[chashu]{bouillon à la viande}}
  \ingr{180}{ml}{ \hyperref[chashu]{marinade du cha-shū}}
  \ingr{1}{\cs}{sauce d'huître}
  \ingr{2}{\cs}{levure maltée}
  \ingr{}{}{sel}
\end{ingredients}

\begin{recipe}
  \begin{enumerate}

  \item Mélanger tous les ingrédients dans une casserole, chauffer et
    ajuster le sel.

  \end{enumerate}
\end{recipe}


\stage{Composition}

\begin{ingredients}
  \ingrS{280--320}{g}{ \hyperref[ラーメン]{nouilles ramen sèches ou fraîches}}
  \ingr{2}{poignées}{pousses de haricots mungo, blanchies 1~min et égouttées}
  \ingr{4}{}{\hyperref[chashu]{ajitsuke tamago}, coupés en deux}
  \ingr{8}{tranches}{ \hyperref[chashu]{cha-shū}}
  \ingr{1}{}{cébette, en très fines rondelles}
  \ingr{\fracH}{feuille}{nori, coupée en 4~rectangles}
  \ingr{}{}{poivre}
  \ingr{}{}{rā-yu (facultatif)}
\end{ingredients}

\begin{recipe}
  \begin{enumerate}

  \item Cuire les nouilles. Égoutter et répartir dans 4~grands bols.

  \item Verser la soupe chaude et disposer la garniture.

  \end{enumerate}
\end{recipe}

%\accord{}
